\section{A uniform freeze-or-grow bound for the external parent}
\label{ps:section}

The statement below applies to the original core process, also allowing
inert boundary stubs.  A stub is counted in its incident core's degree,
receives outward selections without changing state, and sends no arrivals.
No monotonicity of expected descendant counts is used.

\begin{lemma}[Reciprocal population under persistent immigration]
\label{ps:reciprocal}
Let $0<\rho\le1$, and let $X$ be the birth--death chain on
$\mathbb Z_{\ge0}$ with birth rate $x+\rho$ at $x$ and death rate one
at every $x\ge1$.  If $X(0)=0$, then
\begin{equation}\label{ps:reciprocal-bound}
 \mathbb E\frac1{X(t)+1}
 \le \left(1+\frac3\rho\right)e^{-\rho t/3}
 \qquad(t\ge0).
\end{equation}
\end{lemma}
\begin{proof}
Put $q_x=\Pr\{\operatorname{Pois}(1)\ge x\}$ and
\[
 f(x)=\frac1{x+1},\qquad
 h(x)=f(x)+\frac3\rho q_x.
\]
The generator $G_0$ with birth rate $x$ and unit death rate satisfies
$G_0q_x=0$ for $x\ge1$, because
\[
 q_{x+1}-q_x=-\frac{e^{-1}}{x!}.
\]
At $x=0$, $G_0q_0=0$ as well.  For $x\ge1$,
\[
 G_0f(x)=-\frac{x-2}{x(x+2)}.
\]
Consequently, for the generator $G$ including immigration,
\[
 Gh(x)=G_0f(x)
       -\frac{\rho}{(x+1)(x+2)}-\frac3{e\,x!}
 \quad(x\ge1),
\]
whereas $Gh(0)=-\rho/2-3/e$.

For $x\ge1$, summing the factorial tail gives
$q_x\le2/(e\,x!)$.  For $x\ge4$ we also have
$G_0f(x)\le-(5/12)f(x)$.  Since $\rho\le1$, these inequalities imply
\[
 Gh(x)+\frac\rho3 h(x)
 \le-\frac1{12}f(x)-\frac1{e\,x!}<0
 \qquad(x\ge4).
\]
At $x=0,1,2,3$, respectively, the exact expressions for
$Gh(x)+(\rho/3)h(x)$ are
\[
 -\frac\rho6+1-\frac3e,\quad
 \frac43-\frac4e,\quad
 \frac\rho{36}+1-\frac7{2e},\quad
 -\frac1{15}+\frac\rho{30}+1-\frac3e.
\]
They are all negative, using $e<3$ and $\rho\le1$.
Thus $Gh\le-(\rho/3)h$ at every state.

The chain is nonexplosive, and $h$ is bounded.  Localization and
Dynkin's formula therefore yield
$\mathbb E h(X(t))\le h(0)e^{-\rho t/3}$.
Since $f\le h$ and $h(0)=1+3/\rho$, this proves the claim.
\end{proof}

Fix a real core vertex $p$ whose initial full degree is at most an
integer $M\ge2$, and set
\[
 \tau_M=\inf\{t:D_p(t)>M\},\qquad
 C_M=1+3M,\qquad \kappa_M=\frac1{3M}.
\]
For real neighbours $v$ of $p$, define
\[
 \lambda_p(t)=\sum_{v\sim p,\ v\text{ real}}\frac1{D_v(t)}.
\]
Thus $\lambda_p$ is the rate at which $p$ receives a new leaf and
its full degree increases.  Its leaf-promotion rate is
$b_p=L_p/D_p$.

\begin{theorem}[Uniform local stabilization]
\label{ps:stabilization}
Uniformly over all finite initial environments and their boundary stubs,
\begin{equation}\label{ps:rate-bound}
 \mathbb E\bigl[\mathbf1_{\{t<\tau_M\}}\lambda_p(t)\bigr]
 \le M C_M e^{-\kappa_Mt}
       \sum_{j=0}^{M-1}\frac{(C_Mt)^j}{j!}.
\end{equation}
In particular, the probability that $p$ has a degree increase or a leaf
promotion at some time in $(T,\tau_M]$ is at most
\begin{align}
 Q_M(T)
 &=\min\left\{1,
 2M C_M\sum_{j=0}^{M-1}\frac{C_M^j}{j!}
       \int_T^\infty e^{-\kappa_Mt}t^j\,dt\right\}
 \label{ps:tail-integral}\\
 &=\min\left\{1,
 2M C_M e^{-\kappa_MT}
 \sum_{j=0}^{M-1}C_M^j
       \sum_{i=0}^j\frac{T^i}{i!\,\kappa_M^{j-i+1}}
 \right\}.
 \label{ps:tail-explicit}
\end{align}
This bound tends to zero exponentially in $T$ for each fixed $M$.
\end{theorem}
\begin{proof}
While $D_p\le M$, implement the rate-one clock at $p$ using $M$
independent slots, each of rate $1/M$.  Give every neighbour a permanent
distinct slot when that neighbour first appears.  At an occupied slot,
choose its assigned neighbour.  At an empty slot, choose uniformly
among the current $D_p$ neighbours.  Every particular neighbour is
then chosen at rate
\[
 \frac1M+\frac{M-D_p}{M D_p}=\frac1{D_p},
\]
so the construction is exact until $\tau_M$.

Let $v$ be a real neighbour present initially, or born as a new
neighbour of $p$ at a stopping time $\sigma<\tau_M$.  Its occupied
slot supplies arrivals at rate $1/M$.  Every such arrival creates a
leaf at $v$; the first one also turns $v$ into a core if it was still
a leaf.  The number of leaf neighbours of $v$ lower-dominates, until
$\tau_M$, the chain of Lemma~\ref{ps:reciprocal} with $\rho=1/M$
and initial state zero.  Use the same leaf-birth Poisson measure,
the occupied slot for immigration, and a death at every ring of
$v$'s own rate-one clock.  Actual leaf losses occur at most once at
such a ring, and all additional arrivals only increase the actual
leaf count.  These three driving processes are independent before
stopping.  They can be continued after stopping to define the full
comparison chain.

The permanent edge $pv$ gives $D_v\ge X+1$ before stopping.  Therefore,
conditionally on the history at $\sigma$, for every $a\ge0$,
\begin{equation}\label{ps:conditional-kernel}
 \mathbb E\left[
   \frac{\mathbf1_{\{\sigma+a<\tau_M\}}}{D_v(\sigma+a)}
   \,\middle|\,\mathcal F_\sigma\right]
 \le C_Me^{-\kappa_Ma}.
\end{equation}
Initial neighbours use $\sigma=0$.  Stubs contribute zero and require
no comparison chain.

Give each new neighbour $w$ of $p$ a genealogical parent: the old
neighbour $v$ whose clock selected $p$ and created $w$.
This is a genealogy of neighbours, separate from the core genealogy.
Its roots are the at most $M$ initial real neighbours.  Each individual's
conditional mean birth intensity, as a function of its age before
stopping, is bounded by the kernel
$k_M(a)=C_Me^{-\kappa_Ma}$ in
\eqref{ps:conditional-kernel}.
Iterated conditional first moments, or the compensator formula followed
by Tonelli's theorem, bound the expected birth intensity contributed
by neighbour generation $j$ by
\[
 M k_M^{*(j+1)}(t)
 =M C_M^{j+1}e^{-\kappa_Mt}\frac{t^j}{j!}.
\]
No independence between neighbours is asserted or needed.
Before $\tau_M$ there are at most $M$ neighbours, so no contributing
individual can have neighbour generation greater than $M-1$.
Summing gives~\eqref{ps:rate-bound}.

Finally, each leaf neighbour contributes one to $\lambda_p$, so
$b_p=L_p/D_p\le L_p\le\lambda_p$.  The total intensity of changes
to $(D_p,L_p,c_p)$ is thus at most $2\lambda_p$.
The compensator bound for the number of such changes after $T$, counting
also the degree increase at $\tau_M$, proves
\eqref{ps:tail-integral}.  Integrating the integer powers gives
\eqref{ps:tail-explicit}.
\end{proof}

